% 张三慧大学物理第四版V2《固体中的电子》第29章章末习题29.1--29.14
% 原书V2扫描PDF物理页452--453，辅导S PDF386--389。
% 思考题29.1--29.13没有混入正式习题。
\section*{第29章\quad 固体中的电子}
\subsection*{章末习题29.1--29.14}
\paragraph{习题 29.1}
金的密度 \(\rho_m=19.3\,\mathrm{g/cm^3}\)，
假设每原子贡献一个自由电子，计算费米能量、
费米速率、费米温度及具有费米能量的电子
德布罗意波长；金的摩尔质量约
\(197\,\mathrm{g/mol}\)。
\paragraph{习题 29.2}
证明在 \(T=0\) 时三维自由电子费米气
单位体积的总能量和每电子平均能量分别是
\(u=3nE_F/5,\ \bar E=3E_F/5\)。
\paragraph{习题 29.3}
在 \(T=0\) 时求自由电子费米气的
平均速率 \(\bar v\) 和方均根速率 \(v_{\rm rms}\)，
用费米速率 \(v_F\) 表示。
\paragraph{习题 29.4}
中子星由非相对论近似费米中子气构成，
典型密度 \(\rho=5\times10^{16}\,\mathrm{kg/m^3}\)。
求其中子的费米能量及费米速率，
并评估非相对论模型是否充分。
\paragraph{习题 29.5}
电子费米气的费米能量5eV。
在何温度下其电子比热等于经典单原子气体
等容比热的10\%？
按低温Sommerfeld电子比热公式估算。
\paragraph{习题 29.6}
极低温下钾的摩尔热容近似
\[
C_m=(2.08\times10^{-3}T+
2.57\times10^{-3}T^3)\,\mathrm{J/(mol\,K)},
\]
电子贡献与 \(T\) 成正比、晶格点阵贡献与 \(T^3\)
成正比。(1) 由电子线性系数求钾的费米能量；
(2) 求电子与点阵贡献相等的温度。
\paragraph{习题 29.7}
银密度 \(10.5\times10^3\,\mathrm{kg/m^3}\)，
室温电阻率 \(1.6\times10^{-8}\,\Omega\,\mathrm m\)。
每原子提供一个自由电子。
(1) 求自由飞行时间 \(\tau\)；
(2) 按 \(T=300\,\mathrm K\) 的经典热速率估计平均自由程；
(3) 按费米速率估计平均自由程；
(4) 估计点阵离子平均间距并比较(2)(3)。
\paragraph{习题 29.8}
在1000K费米--狄拉克分布下，
相对于费米能量分别高0.10eV和低0.10eV
的量子态，其电子平均占据数各为多少？
\paragraph{习题 29.9}
金刚石禁带宽度5.5eV。
在300K下估计禁带底部与顶部电子占据数之比，
并求使价带电子跨过此禁带所需光子的最大波长。
\paragraph{习题 29.10}
纯硅本征自由电子浓度
\(n_0\simeq10^{16}\,\mathrm{m^{-3}}\)。
通过掺杂磷使电子浓度提高 \(10^6\) 倍，
硅密度 \(2.33\,\mathrm{g/cm^3}\)。
(1) 需以磷替换多少比例的硅原子？
(2) 1.0g硅需增加多少质量的磷？
忽略补偿和不完全电离。
\paragraph{习题 29.11}
硅禁带1.2eV，掺适量磷后施主能级
离导带底仅0.045eV。求此掺杂半导体可吸收
光子的最大波长。
\paragraph{习题 29.12}
CdS和PbS的禁带分别为2.42eV和0.30eV，
求它们作为光电导体时吸收阈值最大波长，
判断所属光谱波段。
\paragraph{习题 29.13}
Ga--As--P发光二极管禁带宽度1.9eV，
求带隙复合光子对应的最大波长。
\paragraph{习题 29.14}
KCl晶体的满价带上方有宽7.6eV禁带。
140nm光照射时，判定在此带间吸收近似下
晶体对该光透明还是不透明。
